% TOP_PROBLEM: {"id": 4800001, "title": "Multiple ergodic averages — Problem 1", "category_id": 11, "category": {"id": 11, "name": "dynamical_systems", "display_name": "Dynamical Systems", "description": "Problems about long-term behavior of deterministic systems, Hamiltonian dynamics, and periodic orbits.", "slug": "dynamical-systems", "order_index": 11, "created_at": "2026-07-31T15:26:25.671Z"}, "status": "open"}
% FIRST_POSTED: null
% ENTRY_KIND: statement_only
% Imported from: batch15.tex; UnsolvedMath ID 4800001
\documentclass[11pt]{article}
\usepackage[margin=25mm]{geometry}
\usepackage{fontspec}
\setmainfont{Latin Modern Roman}
\usepackage{amsmath,amssymb,mathtools}
\usepackage{unicode-math}
\setmathfont{Latin Modern Math}
\usepackage{xurl,hyperref}
\hypersetup{colorlinks=true,urlcolor=blue}
\setcounter{secnumdepth}{0}
\setlength{\parindent}{0pt}
\setlength{\parskip}{5pt}
\setlength{\emergencystretch}{3em}
\newcommand{\sref}[2]{\hyperlink{source:#1:#2}{[S#2]}}
\newcommand{\cataloguescope}[1]{\paragraph{Recorded target.}#1}
\begin{document}
% UnsolvedMath ID 4800001; source catalogue code AMR-047-0001
\setcounter{section}{4800000}
\section{Multiple ergodic averages — Problem 1}\label{4800001}
{\small\sffamily UnsolvedMath ID 4800001 · Dynamical Systems}
\subsection{Definitions and mathematical statement}

\paragraph{Shared notation.}$\mathbb N=\{1,2,\ldots\}$, and $\mathbb Z,\mathbb Q,\mathbb R,\mathbb C$ denote the integers, rationals, reals and complex numbers. Any other conventions are specified locally.


A system is a probability space $(X,\mathcal X,\mu)$ with invertible measure-preserving transformations; multiple transformations are assumed to commute. Write $T^j f=f\circ T^j$, and $T^j A$ for the image of a measurable set under $T^j$. A system with one transformation is ergodic if every invariant measurable set has measure zero or one. The space $L^\infty(\mu)$ consists of essentially bounded complex measurable functions. A nilmanifold of step at most $s$ is $Y=G/\Gamma$, where $G$ is a nilpotent Lie group of step at most $s$ and $\Gamma$ is a discrete cocompact subgroup. A nilsystem is a translation $y\mapsto by$ on $Y$ with its normalized invariant Haar measure. A basic $s$-step nilsequence in $q$ variables is $\psi(n_1,\ldots,n_q)=F(b_1^{n_1}\cdots b_q^{n_q}y)$, with commuting $b_i\in G$, $y\in Y$, and continuous $F:Y\to\mathbb C$. A basic generalized nilsequence permits a bounded Riemann-integrable $F$ instead. A (generalized) nilsequence is a uniform limit of basic (generalized) nilsequences of the same step and number of variables. For $\ell\ge1$ and $f_0,\ldots,f_\ell\in L^\infty(\mu)$, define the bounded multiple-correlation sequence on $\mathbb Z^\ell$ by \[\mathcal C(n_1,\ldots,n_\ell)=\int_X f_0\prod_{j=1}^\ell T_j^{n_j}f_j\,d\mu.\] An integral combination of generalized $\ell$-step nilsequences is a sequence $b(\mathbf n)=\int_Y\psi_y(\mathbf n)\,d\nu(y)$, where $Y$ is a compact metric space, $\nu$ a complex Borel measure of finite total variation, each $\psi_y$ is such a generalized nilsequence and $y\mapsto\psi_y(\mathbf n)$ is measurable and essentially bounded with respect to $|\nu|$ for every $\mathbf n$. An approximate integral combination means that for every $\varepsilon>0$ there is such a $b$ with $\sup_{\mathbf n\in\mathbb Z^\ell}|\mathcal C(\mathbf n)-b(\mathbf n)|\le\varepsilon$.
\paragraph{Statement recorded in the supplied source.}
Determine the structure of all these multiple-correlation sequences. In particular, is every $\mathcal C$ an integral combination, or at least an approximate integral combination, of generalized $\ell$-step nilsequences in $\ell$ variables? \sref{4800001}{2}
\subsection{Short English statement}
Describe every commuting multiple-correlation sequence, and decide whether generalized nilsequences give an exact integral representation or arbitrarily accurate uniform integral approximations.
\subsection{Sources}
\begin{description}
\item[\hypertarget{source:4800001:1}{S1}] ULAM.AI and the UnsolvedMath Contributors, \emph{UnsolvedMath: A Curated Collection of Open Mathematics Problems}, record 4800001 (AMR-047-0001). CC BY 4.0. \url{https://huggingface.co/datasets/ulamai/UnsolvedMath}.
\item[\hypertarget{source:4800001:2}{S2}] N. Frantzikinakis, Some open problems on multiple ergodic averages (2016), Sections 1.2,2.4 and4.1, equation (14), integral-combination definitions and full Problem 1, pp.2,5--6,18. \url{https://arxiv.org/pdf/1103.3808v3}.
\end{description}
\label{end:4800001}
\end{document}
